\setchoices{label=(\Alph*)}

\begin{question}[index=1]
设函数 $f(x)$ 有二阶连续导数，
$\displaystyle \lim_{x\to 0}\frac{f(x)}{\displaystyle\int_0^{x^2}\frac{\sin t}{t}\dd t}=-1$，则
\begin{choices}[columns=1]
\item $f(x)$ 在点 $x=0$ 处取得极大值．
\item $f(x)$ 在点 $x=0$ 处取得极小值．
\item 点 $(0,f(0))$ 是曲线 $y=f(x)$ 的拐点．
\item $x=0$ 不是 $f(x)$ 的极值点，$(0,f(0))$ 也不是曲线 $y=f(x)$ 的拐点．
\end{choices}
\end{question}
\QuestionBreak

\begin{question}[index=2]
设
\[
\begin{aligned}
M&=\int_{-\frac{\pi}{2}}^{\frac{\pi}{2}}\left(\frac{\sin^3 x}{1+x^2}\cos^4 x+1\right)\dd x,\\
N&=\int_{-\frac{\pi}{2}}^{\frac{\pi}{2}}\left(\ln\frac{4+x}{4-x}+\cos^3 x\right)\dd x,\qquad
P=\int_{-\frac{\pi}{2}}^{\frac{\pi}{2}}(\mathrm e^{x^2}\sin x-2)\dd x,
\end{aligned}
\]
则
\begin{choices}[columns=2]
\item $M<N<P$．
\item $M<P<N$．
\item $P<N<M$．
\item $N<P<M$．
\end{choices}
\end{question}
\QuestionBreak

\begin{question}[index=3]
设
\[
f(x,y)=\begin{cases}
\dfrac{xy}{x^2+y^2}, & x^2+y^2\ne0,\\[3pt]
0, & x^2+y^2=0,
\end{cases}
\]
则函数 $f(x,y)$
\begin{choices}[columns=2]
\item 在点 $(0,0)$ 处连续．
\item 在点 $(0,0)$ 处一阶偏导数不存在．
\item 在点 $(0,0)$ 的邻域内无界．
\item 与 $\sin(x+y)$ 之积在点 $(0,0)$ 处连续．
\end{choices}
\end{question}
\QuestionBreak

\begin{question}[index=4]
设空间曲面 $\Sigma_1:z=\dfrac{x^3+y^3}{3}$，$\Sigma_2:z=\dfrac{x^2-y^2}{2}$，
$\Sigma_3:z=\dfrac{x^2y^2}{2}$ 被柱面 $x^2+y^2=1$ 所截部分的面积分别为 $S_1,S_2,S_3$，则
\begin{choices}[columns=2]
\item $S_1>S_2>S_3$．
\item $S_2>S_1>S_3$．
\item $S_3>S_1>S_2$．
\item $S_3>S_2>S_1$．
\end{choices}
\end{question}
\QuestionBreak

\begin{question}[index=5]
设 $\symbfit{A}$ 是 $3\times2$ 矩阵，且 $\symbfit{A}^{\mathrm T}\symbfit{A}=\symbfit{E}$，
$\symbfit{\alpha}$ 为 $3$ 维非零列向量，$\symbfit{\beta}=\symbfit{A}\symbfit{A}^{\mathrm T}\symbfit{\alpha}$，则以下结论：
\par ① $\symbfit{A}$ 的列向量正交；
\par ② $\symbfit{A}\symbfit{x}=\symbfit{\beta}$ 无解；
\par ③ $\symbfit{A}^{\mathrm T}\symbfit{\alpha}=\symbfit{A}^{\mathrm T}\symbfit{\beta}$；
\par ④ $\symbfit{\beta}$ 与 $\symbfit{\alpha}-\symbfit{\beta}$ 正交．
\par 所有正确结论的序号为
\begin{choices}[columns=2]
\item ①③．
\item ②④．
\item ①②③．
\item ①③④．
\end{choices}
\end{question}
\QuestionBreak

\begin{question}[index=6]
设 $\symbfit{A}$ 是 $n$ 阶实对称矩阵，则“存在 $n$ 阶矩阵 $\symbfit{B}$，使
$\symbfit{A}\symbfit{B}+\symbfit{B}^{\mathrm T}\symbfit{A}$ 是正定矩阵”是“$r(\symbfit{A})=n$”的
\begin{choices}[columns=2]
\item 充分非必要条件．
\item 必要非充分条件．
\item 充分必要条件．
\item 既非充分也非必要条件．
\end{choices}
\end{question}
\QuestionBreak

\begin{question}[index=7]
设 $f(x_1,x_2,x_3)=\symbfit{x}^{\mathrm T}\symbfit{A}\symbfit{x}$，$\symbfit{A}=\symbfit{A}^{\mathrm T}$，
$\symbfit{x}=(x_1,x_2,x_3)^{\mathrm T}$，$\symbfit{A}$ 的特征值为 $\lambda_1,\lambda_2,\lambda_3$，
且 $\lambda_1>\lambda_2>\lambda_3$，对应的特征向量分别为 $\symbfit{\xi}_1,\symbfit{\xi}_2,\symbfit{\xi}_3$．
若 $\symbfit{x}^{\mathrm T}\symbfit{x}=1$，$\symbfit{x}^{\mathrm T}\symbfit{\xi}_1=0$，则 $f(x_1,x_2,x_3)$ 的最大值为
\begin{choices}[columns=4]
\item $\lambda_1$．
\item $\lambda_2$．
\item $\lambda_1+\lambda_2$．
\item $\lambda_2+\lambda_3$．
\end{choices}
\end{question}
\QuestionBreak

\begin{question}[index=8]
设 $(X,Y)$ 服从二维正态分布，$EX=\mu_1$，$DX=\sigma_1^2$，$EY=\mu_2$，$DY=\sigma_2^2$，
$\rho_{XY}=\rho$，且 $\mu_1\ne\mu_2$，$\sigma_1\ne\sigma_2$．记 $\xi=aX+bY$，$\eta=aX-bY$，
$ab\ne0$．如果 $\xi$ 与 $\eta$ 相互独立，则 $a$ 与 $b$ 应满足关系式
\begin{choices}[columns=2,label-pos=left]
\item $\dfrac{a^2}{b^2}=\dfrac{\sigma_2^2}{\sigma_1^2}$．
\item $\dfrac{a^2}{b^2}=\dfrac{\sigma_1^2}{\sigma_2^2}$．
\item $\dfrac{a^2}{b^2}=\dfrac{\mu_2^2}{\mu_1^2}$．
\item $\dfrac{a^2}{b^2}=\dfrac{\mu_1^2}{\mu_2^2}$．
\end{choices}
\end{question}
\QuestionBreak

\begin{question}[index=9]
设随机事件 $A,B$ 满足 $0<P(A)<1$，$0<P(B)<1$，则 $P(AB)>P(A)P(B)$ 的充要条件是
\begin{choices}[columns=2]
\item $P(A\overline B)>P(A)P(\overline B)$．
\item $P(\overline A\overline B)>P(\overline A)P(\overline B)$．
\item $P(B\mid\overline A)>P(B\mid A)$．
\item $P(A\mid\overline B)>P(A\mid B)$．
\end{choices}
\end{question}
\QuestionBreak

\begin{question}[index=10]
设随机变量 $X,Y$ 独立同分布于 $E(\lambda)$，其中 $\lambda>0$，$F(x)$ 为 $X$ 的分布函数，
则下列随机变量中与 $F(X)$ 同分布的是
\begin{choices}[columns=4,label-pos=left]
\item $\dfrac{2X}{X+Y}$．
\item $\dfrac{X}{Y}$．
\item $\dfrac{X+Y}{2X}$．
\item $\dfrac{Y}{X+Y}$．
\end{choices}
\end{question}
